% Theory-only supplemental material for the Rayleigh-BIC manuscript.
\documentclass[aps,prl,reprint,superscriptaddress,nofootinbib,floatfix]{revtex4-2}

\usepackage[T1]{fontenc}
\usepackage{lmodern}
\usepackage{microtype}
\usepackage{amsmath,amssymb,bm}
\usepackage{graphicx}
\usepackage{xcolor}
\usepackage{booktabs}
\usepackage[hidelinks]{hyperref}

\hypersetup{
  pdftitle={Supplemental theory for a strict Rayleigh BIC in an underwater acoustic metagrating},
  pdfauthor={Author names to be inserted}
}

\newcommand{\dd}{\,\mathrm{d}}
\newcommand{\ii}{\mathrm{i}}
\newcommand{\e}{\mathrm{e}}
\newcommand{\vect}[1]{\bm{#1}}
\newcommand{\RA}{\mathrm{RA}}
\providecommand{\dr}{\mathrm{dr}}
\providecommand{\hom}{\mathrm{hom}}
\providecommand{\inc}{\mathrm{inc}}
\newcommand{\figplaceholder}[2]{%
  \IfFileExists{#1}{\includegraphics[width=\linewidth]{#1}}{%
    \fbox{\begin{minipage}[c][0.24\textheight][c]{0.92\linewidth}
      \centering\small
      \textit{Vector supplemental figure supplied separately}\par\vspace{0.5ex}
      \texttt{\detokenize{#1}}\par\vspace{0.8ex}
      #2
    \end{minipage}}}}

\renewcommand{\theequation}{S\arabic{equation}}
\renewcommand{\thetable}{S\arabic{table}}
\renewcommand{\thefigure}{S\arabic{figure}}

\begin{document}
\title{Supplemental theory: a strict Rayleigh bound state in the continuum in a two-groove underwater acoustic metagrating}
\author{Author Name}
\affiliation{Affiliation to be inserted}
\author{Coauthor Name}
\affiliation{Affiliation to be inserted}
\date{\today}
\maketitle

\section{Scope and notation}
\label{sec:s-scope}

This supplement gives the modal derivation and numerical diagnostics behind
the theory-only main text.  It does not add an experiment, a finite-element
calculation, or a loss model.  The water sound speed is $c=1500$~m/s, the
nominal frequency is $f_0=200$~kHz, and the time convention is
$\e^{+\ii\omega t}$.  Superscripts $\dr$ and $\hom$ distinguish driven
scattering amplitudes from outgoing amplitudes of a homogeneous eigenmode.
The strict Rayleigh criterion is imposed on the latter.

\section{Floquet expansion and groove matching}
\label{sec:s-floquet}

For a period $a$ and incidence angle $\theta$, define
\begin{equation}
 k_0=\frac{\omega}{c},\qquad k_{x,\inc}=k_0\sin\theta,
 \qquad k_{y,\inc}=k_0\cos\theta .
 \label{eq:s-wavenumbers}
\end{equation}
The incident field travels toward $-y$ and the outgoing reflected expansion
above the surface is
\begin{equation}
 p_{\inc}=\e^{\ii(k_{x,\inc}x+k_{y,\inc}y)},
 \qquad
 p_{\mathrm{sc}}=\sum_n A_n^{\dr}\e^{\ii k_{x,n}x-\ii k_{y,n}y},
 \label{eq:s-exterior}
\end{equation}
\begin{equation}
\begin{aligned}
 k_{x,n}&=k_{x,\inc}+\frac{2\pi n}{a},\\
 k_{y,n}&=\sqrt{k_0^2-k_{x,n}^2},\\
 \Re k_{y,n}&\geq0,\qquad \Im k_{y,n}\leq0 .
\end{aligned}
 \label{eq:s-floquet}
\end{equation}
With $\Omega=a/\lambda$ and $\kappa=a k_{x,\inc}/(2\pi)$, the Rayleigh
condition for order $n$ is
\begin{equation}
 |\kappa+n|=\Omega .
 \label{eq:s-ra}
\end{equation}

For groove $\ell$, let $x_\ell$ be the left edge and let the groove occupy
$x_\ell<x<x_\ell+w_\ell$ and $-d_\ell<y<0$.  A bottom-rigid cosine expansion
is
\begin{equation}
\begin{aligned}
 p_\ell(x,y)&=\sum_{q=0}^{K-1} C_{\ell q}
 \cos\!\left[\frac{q\pi(x-x_\ell)}{w_\ell}\right]\\[-2pt]
 &\quad\times\cos\!\left[\beta_{\ell q}(y+d_\ell)\right],\\
 \beta_{\ell q}^2&=k_0^2-\left(\frac{q\pi}{w_\ell}\right)^2 .
\end{aligned}
 \label{eq:s-groove}
\end{equation}
The aperture projection used by the matching calculation is
\begin{equation}
 \Pi_{n,\ell q}=\frac{1}{w_\ell}\int_{x_\ell}^{x_\ell+w_\ell}
 \e^{-\ii k_{x,n}x}
 \cos\!\left[\frac{q\pi(x-x_\ell)}{w_\ell}\right]\dd x .
 \label{eq:s-projection}
\end{equation}
The cosine norm is $1$ for $q=0$ and $1/2$ for $q>0$; the corresponding
factors are included in the numerical projection matrices.

For the chosen time convention, pressure-normalized normal velocity is
$\rho\omega v_y=+\ii\partial_y p$.  At the groove aperture this gives
\begin{equation}
 \rho\omega v_{y,\ell}(x,0)=
 -\ii\sum_q \beta_{\ell q} C_{\ell q}
 \sin(\beta_{\ell q}d_\ell)
 \cos\!\left[\frac{q\pi(x-x_\ell)}{w_\ell}\right].
 \label{eq:s-groove-velocity}
\end{equation}
The exterior reflected contribution has
$\rho\omega v_{y,n}=k_{y,n}A_n^{\dr}$, while the incident contribution has
$\rho\omega v_{y,\inc}=-k_{y,\inc}p_{\inc}$.

\subsection{Block operator}
\label{sec:s-block}

Let $\mathsf B$ project external pressure onto aperture cosine modes and let
$\mathsf D$ project aperture cosine velocities onto global Floquet tests.  If
$\mathsf C_d=\operatorname{diag}[\cos(\beta_{\ell q}d_\ell)]$ and
$\mathsf V_d=\operatorname{diag}[\beta_{\ell q}\sin(\beta_{\ell q}d_\ell)]$, the
homogeneous pressure--velocity system can be written
\begin{equation}
 \mathsf F
 \begin{bmatrix}\vect{A}^{\hom}\\ \vect{C}\end{bmatrix}=0,
 \qquad
 \mathsf F=
 \begin{bmatrix}
  \mathsf B & -\mathsf C_d\\[2pt]
  \operatorname{diag}(k_{y,n}/k_0) & \ii\mathsf D\mathsf V_d
 \end{bmatrix} .
 \label{eq:s-block}
\end{equation}
The top rows enforce pressure continuity in each aperture.  The lower rows
enforce normal-velocity continuity after the rigid portions of the surface
are included as zero-velocity regions.  Keeping $C_{\ell q}$ as independent
unknowns avoids the tangent poles that appear if the groove field is first
eliminated in a driven calculation.

For the forced problem, the same continuity equations can be reduced to
\begin{equation}
\begin{aligned}
 \left[\operatorname{diag}(k_{y,n})-\mathsf S\mathsf B\right]
 \vect{A}^{\dr}&=k_{y,\inc}\vect e_0+\mathsf S\vect b_{\inc},\\
 \mathsf S&=\mathsf D\operatorname{diag}
 \left[-\ii\beta_{\ell q}\tan(\beta_{\ell q}d_\ell)\right].
\end{aligned}
 \label{eq:s-driven}
\end{equation}
The vector in Eq.~\eqref{eq:s-driven} contains the incident aperture
projections.  A driven $A_0^{\dr}$ therefore includes the background
reflection and is not the homogeneous $A_0^{\hom}$ used below.

\section{Strict zero amplitudes and square integrability}
\label{sec:s-strict}

The normal-power efficiency of a propagating driven order is
\begin{equation}
 \eta_n=\frac{\Re k_{y,n}}{k_{y,\inc}}
 \left|\frac{A_n^{\dr}}{A_{\inc}}\right|^2,
 \qquad A_{\inc}=1 .
 \label{eq:s-efficiency}
\end{equation}
This expression alone cannot certify a BIC at a threshold.  To see why, take
one exterior term $p_n=A_n^{\hom}\e^{\ii k_{x,n}x-\ii k_{y,n}y}$ and integrate its
pressure norm from the surface to height $H$.  For a threshold order,
$k_{y,n}=0$ and
\begin{equation}
 \int_0^H |p_n|^2\dd y=H|A_n^{\hom}|^2 .
 \label{eq:s-threshold-norm}
\end{equation}
The same linear growth occurs for a propagating order after averaging over
oscillations.  For an evanescent order with $k_{y,n}=-\ii\alpha_n$ and
$\alpha_n>0$, the integral saturates as
\begin{equation}
 \int_0^\infty |p_n|^2\dd y=\frac{|A_n^{\hom}|^2}{2\alpha_n} .
 \label{eq:s-evanescent-norm}
\end{equation}
Consequently, at the positive-angle design point the strict condition is
\begin{equation}
 A_{-1}^{\hom}=0,\qquad A_0^{\hom}=0,
 \label{eq:s-strict-positive}
\end{equation}
where $n=-1$ is the selected RA order and $n=0$ is the finite-flux open order.
At the time-reversed point the target is $n=+1$ and the second condition
remains $A_0^{\hom}=0$.  Zero grazing flux without
$A_{n_{\RA}}^{\hom}=0$ is therefore insufficient.

\subsection{Radiation-reduced SVD}
\label{sec:s-svd}

Let $\mathsf F_{\mathrm{red}}$ be obtained from $\mathsf F$ by removing the
columns corresponding to the threshold amplitude and all finite-flux open
amplitudes.  The remaining groove columns are retained.  To avoid a spurious
minimum caused by disparate modal scales, define
\begin{equation}
\begin{aligned}
 \mathsf R&=\operatorname{diag}(r_i),&
 r_i&=\max_j|F_{\mathrm{red},ij}|_2,\\
 \mathsf C&=\operatorname{diag}(c_j),&
 c_j&=\max_i|[\mathsf R^{-1}\mathsf F_{\mathrm{red}}]_{ij}|_2,
\end{aligned}
 \label{eq:s-scaling}
\end{equation}
and
\begin{equation}
 \widetilde{\mathsf F}=\mathsf R^{-1}\mathsf F_{\mathrm{red}}\mathsf C^{-1}.
 \label{eq:s-scaled}
\end{equation}
If $\vect y$ is the right singular vector associated with the smallest
singular value of $\widetilde{\mathsf F}$, the physical reduced vector is
$\vect u=\mathsf C^{-1}\vect y$.  The deleted amplitudes are reinserted as
exact zeros to form $\vect A^{\hom}$.  We report both
\begin{equation}
 \sigma_{\mathrm{rel}}=\frac{\sigma_{\min}}{\sigma_{\max}},
 \qquad
 r_{\mathrm{raw}}=\frac{\|\mathsf F_{\mathrm{red}}\vect u\|_2}{\|\vect u\|_2} .
 \label{eq:s-residuals}
\end{equation}
At $(N,K)=(313,39)$ these are $1.2604\times10^{-15}$ and
$1.3766\times10^{-13}$, respectively.  The two residuals probe different
coordinates and should not be interpreted as two independent physical
observables.

\begin{table*}[t]
\caption{\label{tab:s-strict-convergence}Strict and driven convergence data retained for the theory draft.  A dash means that a quantity was not used as the acceptance diagnostic at that truncation.}
\begin{ruledtabular}
\begin{tabular}{cccccc}
$(N,K)$ & $\sigma_{\mathrm{rel}}$ & $r_{\mathrm{raw}}$ & $f_{\mathrm{router}}$ (kHz) & $\eta_{\mathrm{anom}}$ & role\\
\hline
$(121,15)$ & --- & --- & 200.21475 & 0.30737 & router retuning\\
$(185,23)$ & --- & --- & --- & 0.30662 & router retuning\\
$(281,35)$ & --- & --- & --- & 0.30620 & router retuning\\
$(313,39)$ & $1.2604\times10^{-15}$ & $1.3766\times10^{-13}$ & 200.21687 & 0.30613 & strict BIC / router\\
\end{tabular}
\end{ruledtabular}
\end{table*}

\begin{table}[t]
\caption{\label{tab:s-off-rayleigh}Same-cell off-Rayleigh anomalous-reflection convergence.  The target is $n=-1$ for positive incidence and $n=+1$ for negative incidence.}
\begin{ruledtabular}
\begin{tabular}{ccc}
$(N,K)$ & $\eta_{\mathrm{target}}$ & $\eta_0$\\
\hline
$(81,11)$ & 0.97791 & 0.02209\\
$(121,15)$ & 0.99298 & 0.00702\\
$(161,20)$ & 0.99864 & 0.00136\\
$(201,25)$ & 0.99963 & 0.00037\\
$(313,39)$ & 0.99992 & 0.00008\\
$(401,50)$ & 0.9996185 & 0.0003815\\
\end{tabular}
\end{ruledtabular}
\end{table}

\section{Pole continuation and its caveats}
\label{sec:s-pole}

The complex pole is continued from the Rayleigh point by specifying the
Riemann sheet of the selected order.  For the positive-angle design the target
is $n=-1$; the non-target open orders remain on the outgoing sheet and closed
orders remain on the decaying sheet.  A local complex coordinate $\bar q$ is
used for the target normal wavenumber, and the pole frequency is reconstructed
from
\begin{equation}
 \Omega_p=\sqrt{(\kappa-1)^2+\bar q^{\,2}},
 \qquad
 Q=\frac{\Re\Omega_p}{2|\Im\Omega_p|} .
 \label{eq:s-pole-q}
\end{equation}
The continuation used for the reported local fit retains $(N,K)=(121,15)$.
It reaches a nearest finite $Q$ of about $2\times10^{12}$ before
double-precision and sheet-tracking errors become material.  The local fit
is approximately quadratic, $Q\propto|\Delta\kappa|^{-2.00}$, over the resolved interval.  This is a
finite-model numerical trend, not a universal exponent and not a theorem
about the continuum problem.  The pole diagnostic is also not mixed with the
$(313,39)$ strict-root residual in Eq.~\eqref{eq:s-residuals}.

\section{Single-cavity ablation and interpretation}
\label{sec:s-ablation}

The two-cavity design was compared with a single large rectangular cavity while
retaining additional transverse cosine modes.  In the searched
rectangular-groove family, the best complete strict residual of the one-cavity
model is approximately $0.124$, whereas the two-cavity model reaches the
residual in Eq.~\eqref{eq:s-residuals}.  The result means that additional
transverse modes in that one cavity did not supply the independent phase needed
to cancel both $A_0^{\hom}$ and $A_{-1}^{\hom}$.  It is not a mathematical
no-go theorem for a different one-cavity shape, a compound cavity, or a model
with an additional internal resonator.

The mechanism can be written as two radiation phasor equations,
\begin{equation}
 A_m^{\hom}=\sum_{\ell,q}\mathcal R_{m,\ell q}C_{\ell q}=0,
 \qquad m\in\{0,-1\},
 \label{eq:s-phasors}
\end{equation}
where the wide cavity mainly supplies stored energy and the smaller cavity
provides an independent phase trim.  The roles are geometric and modal, not
experimental assignments.

\section{Dimensional and environmental sensitivity}
\label{sec:s-tolerance}

The design scales with $a=\Omega c/f$.  If the sound speed changes from
$1500$~m/s to $1480$~m/s while the geometry and nominal frequency are held
fixed, the same normalized feature is expected near
$f\simeq197.33$~kHz.  This scaling is a guide for retuning, not a measured
temperature calibration.  A complete tolerance study should perturb
$d_1,d_2,w_1,w_2,$ and $g$ independently, recompute the strict residual and
the driven efficiency, and report the matrix condition number alongside the
efficiency.  The near-BIC point at $\theta=5.27044^{\circ}$ is especially
ill-conditioned ($\operatorname{cond}\simeq1.6\times10^{10}$), so fixed-geometry
efficiency is not a sufficient robustness metric there.

The off-Rayleigh operating point at $\theta=\pm32.6328^{\circ}$ and
$f=202.430$~kHz is separated from this sensitivity discussion.  Its active
$n=-1$ threshold is approximately $141.825$~kHz, so the $99.96185\%$ target
efficiency is not a Rayleigh-critical or BIC-enhanced value.

\section{Energy, reciprocity, and reproducibility checks}
\label{sec:s-checks}

The Floquet normalization in Eq.~\eqref{eq:s-efficiency} makes the
propagating-order basis power normalized.  We use the ordered channel bases
$(0,-1)$ at positive incidence and $(0,+1)$ at negative incidence.  In the lossless model the driven
calculation obeys
\begin{equation}
 \delta_E=\left|1-\sum_{n\in\mathrm{prop}}\eta_n\right|\ll1 ,
 \label{eq:s-energy-error}
\end{equation}
and the power-normalized two-channel matrix satisfies
\begin{equation}
\begin{aligned}
 \mathsf S_\pm^\dagger\mathsf S_\pm&=\mathsf I,\\
 \mathsf S_+(\kappa)&=\mathsf S_-^{\mathsf T}(-\kappa),\\
 \eta_{-1}(+\theta)&=\eta_{+1}(-\theta).
\end{aligned}
 \label{eq:s-reciprocity}
\end{equation}
At the off-Rayleigh point, $\delta_E=4.4\times10^{-16}$ at $(N,K)=(401,50)$.
This is a consistency check of the lossless driven solver, not evidence that
losses vanish in a future finite sample.

At the positive-angle Rayleigh point, $k_{y,-1}=0$.  If
$A_{-1}^{\hom}\neq0$, the threshold pressure is independent of $y$ and its
height-integrated norm grows linearly, as shown in Eq.~\eqref{eq:s-threshold-norm}.
The strict zero in Eq.~\eqref{eq:s-strict-positive} removes this contribution;
the retained $n=+1$ and higher orders are evanescent and have finite integrated
norm.  Complex conjugation and order reversal give the negative-angle state.

The nominal dimensional inputs are $f_0=200\times10^3$~Hz and $c=1500$~m/s.
The normalized design is converted using $a=\Omega c/f_0$ and
$\sin\theta=\kappa/\Omega$.  The external expansion retains an odd number
$N$ of orders centered on $n=0$, and each groove retains $K$ cosine modes.
The strict result uses $(N,K)=(313,39)$; the independent off-Rayleigh driven
point uses $(401,50)$.  Before the singular-value decomposition, the strict
operator is row- and column-scaled.  The threshold and finite-flux columns
are removed, the smallest right singular vector is reconstructed in physical
coordinates, and both the scaled ratio and raw unscaled residual are recorded.
Pole continuation uses the selected outgoing/decaying Riemann sheet.  Driven
efficiencies and energy closure are evaluated over all propagating orders
retained by the external truncation.

\section{PROPOSED future water-tank experiment}
\label{sec:s-experiment}

This section is a proposed protocol for a later paper.  It records planned
controls and observables only; no experiment has been performed for the
present theory draft.

\begin{figure*}[t]
  \centering
  \figplaceholder{figures/figS1_experiment_protocol.pdf}{Proposed future measurement and data-reduction protocol.}
  \caption{\label{fig:s-protocol}
  PROPOSED future water-tank protocol, shown for planning only.  (a) Tank
  geometry and the calibrated flat-plate reference.  (b) The proposed
  angle--frequency scan sequence.  (c) Complex near-field sampling followed
  by Floquet-amplitude extraction.  (d) The proposed ring-down check.  All
  panels are schematics or theory plans; no measurement is reported.}
\end{figure*}

\begin{figure*}[t]
  \centering
  \figplaceholder{figures/figS2_controls_tolerance.pdf}{Control samples, dimensional tolerance, and sound-speed sensitivity plan.}
  \caption{\label{fig:s-controls}
  PROPOSED controls and sensitivity plan.  (a) Full, single-wide-cavity, and
  blocked-small-cavity controls.  (b) The finite-array plan with $M=20$, $40$,
  and $60$ periods.  (c) The cached theory-only dimensional-tolerance
  prediction.  (d) The theory-only sound-speed scaling prediction.  No
  control measurement or completed tolerance map is claimed.}
\end{figure*}

\subsection{Samples and calibration}
\label{sec:s-samples}

The proposed sample set contains an infinite-period surrogate in simulation,
finite two-groove arrays with $M=20$, $40$, and $60$ periods, a flat rigid
reference plate, a single-wide-cavity control, and a full sample with the
small cavity blocked.  The material, backing thickness, and fluid--structure
compliance must be measured or modeled before calling a sample rigid.  The
water temperature and sound speed should be recorded by time of flight.  The
source angular spectrum, hydrophone transfer function, and flat-plate complex
reflection should be calibrated before scanning the grating.

\subsection{Scan and field extraction}
\label{sec:s-scan}

The proposed scan first maps approximately $190$--$205$~kHz over
$-7^{\circ}\leq\theta\leq7^{\circ}$, then refines the features near
$\pm5.25^{\circ}$ and $\pm6^{\circ}$.  An independent map near
$\pm32.63^{\circ}$ should test the off-Rayleigh point.  Complex pressure can
be sampled over several periods at a set height $y_0$ and fitted to
\begin{equation}
\begin{aligned}
 p_{\mathrm{sc}}(x,y_0)&=\sum_n \widetilde A_n(y_0)\e^{\ii k_{x,n}x},\\
 A_n&=\widetilde A_n(y_0)\e^{+\ii k_{y,n}y_0},
\end{aligned}
 \label{eq:s-field-fit}
\end{equation}
then inserted into Eq.~\eqref{eq:s-efficiency}.  The fit should use a measured
incident field and a flat-plate reference so that the driven $A_0^{\dr}$ is
not confused with $A_0^{\hom}$.  Near-field scans and far-field angular scans
should be analyzed together, because a grazing pressure feature can carry
little normal flux.

\subsection{Ring-down and falsifiable checks}
\label{sec:s-ringdown}

For a narrowband burst, switching off the source and fitting an amplitude decay
$\exp(-t/\tau_A)$ would give the operational estimate
\begin{equation}
 Q_{\mathrm{ring}}=\pi f\tau_A .
 \label{eq:s-ringdown}
\end{equation}
The decisive test would be the joint convergence of the extracted complex
Floquet amplitudes, the threshold-order pressure, the finite-flux power, and
the reciprocal signed-order pair.  A driven reflection dip alone would not
validate the BIC.  Finite width, source aperture, disorder, viscous loss, and
fluid--structure motion must be included in the uncertainty budget.

\section{Data and code statement}
\label{sec:s-code}

The theory values are reproducible from the modal-matching implementation in
the repository's `Ni2019\_MATLAB/` directory.  The relevant routines construct
the driven modal solver, the pole-free full homogeneous operator, the strict
Rayleigh reduced operator, the complex-pole continuation, the reciprocal
direction pair, and the routing objective.  The present supplement records the
equations and numerical settings needed to reimplement those calculations; it
does not depend on a finite-element package.  All proposed experiment figures
in Figs.~\ref{fig:s-protocol} and \ref{fig:s-controls} are marked as proposed
and must not be read as completed measurements.

\end{document}
